\section{Minimax rates and consistency}

The minimax squared-error risk \leanref{sym:\mathfrak R_{n,d,\epsilon}}{\ensuremath{\mathfrak R_{n,d,\epsilon}}}, defined in \cref{obj:def:minimax-risk}, measures the precision attainable uniformly over the observed class. Its rate depends jointly on the number of covariate cells and the sample-size scale associated with the public overlap floor. The following result characterizes this dependence with universal comparison constants and identifies a deterministic estimator attaining the rate.

\begin{theoremv}{thm:matched-frontier}[Universal matched risk frontier]
There exist universal real constants \(c_\star,C_\star,H_0,\kappa\) and an integer \(D_0\) satisfying
\[
0<c_\star\le C_\star,\qquad H_0>0,\qquad \kappa>0,\qquad D_0\ge2,
\]
such that the following bounds hold for every:
\begin{itemize}
\item \textbf{(Sample size.)} Integer \(n\ge1\).
\item \textbf{(Alphabet size.)} Integer \(d\ge2\).
\item \textbf{(Overlap.)} Real overlap parameter \(0<\epsilon\le1/2\).
\end{itemize}
Write \(L_d=\log(ed)\) for the logarithmic alphabet scale, and let
\(\widehat V_{n,d,\epsilon}^{\mathrm{AF}}\) be the prescribed deterministic armwise estimator with calibration constants \(H_0,\kappa,D_0\). The observed minimax squared-error risk,
\[
\mathfrak R_{n,d,\epsilon}
=
\inf_{\widehat V}
\sup_{\mathbb P\in\mathcal D_{d,\epsilon}^{\mathrm{obs}}}
E_{\mathbb P}\!\left[
\bigl(\widehat V-\Psi(\mathbb P)\bigr)^2
\right],
\]
where the infimum ranges over measurable real-valued estimators based on \(n\) observed units, satisfies
\[
c_\star\min\left\{1,\frac{d}{n\epsilon L_d}\right\}
\le
\mathfrak R_{n,d,\epsilon}
\le
\sup_{\mathbb P\in\mathcal D_{d,\epsilon}^{\mathrm{obs}}}
E_{\mathbb P}\!\left[
\bigl(\widehat V_{n,d,\epsilon}^{\mathrm{AF}}-\Psi(\mathbb P)\bigr)^2
\right]
\le
C_\star\min\left\{1,\frac{d}{n\epsilon L_d}\right\}.
\]
For the causal class, the same estimator and constants satisfy
\[
\begin{aligned}
c_\star\min\left\{1,\frac{d}{n\epsilon L_d}\right\}
&\le
\inf_{\widehat V}
\sup_{P\in\mathcal P_{d,\epsilon}}
E_{\operatorname{Obs}(P)}\!\left[
\bigl(\widehat V-V^\star(P)\bigr)^2
\right]
\\
&\le
\sup_{P\in\mathcal P_{d,\epsilon}}
E_{\operatorname{Obs}(P)}\!\left[
\bigl(\widehat V_{n,d,\epsilon}^{\mathrm{AF}}-V^\star(P)\bigr)^2
\right]
\\
&\le
C_\star\min\left\{1,\frac{d}{n\epsilon L_d}\right\},
\end{aligned}
\]
where the infimum again ranges over measurable real-valued estimators based on \(n\) observed units, \(V^\star(P)\) is the causal oracle value, and the expectation is under the observed marginal law \(\operatorname{Obs}(P)\).
\end{theoremv}

The ratio \(d/(n\epsilon L_d)\) summarizes the joint statistical cost of many cells and weak overlap. Here \(n\epsilon\) is the rare-arm sample-size scale permitted by the overlap floor: when an arm's treatment probability equals that floor throughout the population, its expected sample count is \(n\epsilon\). Holding \(n\) and \(\epsilon\) fixed, the characterized comparison scale increases with \(d\) through \(d/\log(ed)\). In the saturation regime \(n\epsilon\le d/L_d\), the minimax risk is bounded above and below by positive universal constants. When \(n\epsilon>d/L_d\), the squared-error rate is \(d/(n\epsilon L_d)\). The causal interpretation of these bounds follows from \cref{obj:prop:identification}, which connects the observed target and class to the causal oracle value and its admissible completions.

Along sequences in which the cell count and overlap floor vary with sample size, the same bounds give an exact criterion for uniform mean-square consistency.

\begin{theoremv}{thm:consistency-threshold}[Consistency threshold and fixed-parameter rates]
Let \(\mathfrak R_{n,d,\epsilon}\) be the minimax squared-error risk in \cref{obj:def:minimax-risk}. Consider sequences \((d_n)_{n\in\mathbb N}\) of alphabet sizes and \((\epsilon_n)_{n\in\mathbb N}\) of overlap floors satisfying:
\begin{itemize}
\item \textbf{(Integer alphabet.)} \(d_n\in\mathbb N\) for every \(n\).
\item \textbf{(Alphabet size.)} \(d_n\ge2\) for every \(n\).
\item \textbf{(Positive overlap.)} \(\epsilon_n>0\) for every \(n\).
\item \textbf{(Overlap range.)} \(\epsilon_n\le1/2\) for every \(n\).
\end{itemize}
There exists a sequence of measurable real-valued estimators \(\widehat V_n\), each based on \(n\) independent observed units, such that
\[
\sup_{\mathbb P\in\mathcal D_{d_n,\epsilon_n}^{\mathrm{obs}}}
\mathbb E_{\mathbb P}
\bigl[(\widehat V_n-\Psi(\mathbb P))^2\bigr]
\longrightarrow0
\]
if and only if
\[
\frac{d_n}{n\epsilon_n\log(e d_n)}
\longrightarrow0
\qquad (n\to\infty).
\]

For each fixed integer \(d\ge2\), there exist real constants \(c,C\), depending on \(d\), with \(0<c\le C\), such that for every integer \(n\ge1\) and every overlap floor \(0<\epsilon\le1/2\),
\[
c\min\left\{1,\frac{1}{n\epsilon}\right\}
\le \mathfrak R_{n,d,\epsilon}
\le C\min\left\{1,\frac{1}{n\epsilon}\right\}.
\]

For each fixed overlap floor \(0<\epsilon\le1/2\), there exist real constants \(c,C\), depending on \(\epsilon\), with \(0<c\le C\), such that for every pair of integers \(n\ge1\) and \(d\ge2\),
\[
c\min\left\{1,\frac{d}{n\log(e d)}\right\}
\le \mathfrak R_{n,d,\epsilon}
\le C\min\left\{1,\frac{d}{n\log(e d)}\right\}.
\]
\end{theoremv}

The criterion requires the rare-arm sample-size scale to grow faster than the cell count divided by its logarithmic scale. It therefore allows growing alphabets and shrinking overlap simultaneously, provided their joint burden vanishes relative to the available sample. For a fixed alphabet, the cell-count factor is absorbed into constants depending on that alphabet, yielding the rate \(\min\{1,1/(n\epsilon)\}\); consistency then requires \(n\epsilon_n\to\infty\). For fixed overlap, constants may depend on the overlap floor, and the rate becomes \(\min\{1,d/[n\log(ed)]\}\). These specializations make the joint criterion directly interpretable for designs with either a stable covariate alphabet or a stable overlap floor.

The attaining construction is specified in \cref{obj:def:armwise-estimator}, and its uniform risk bound is established in \cref{obj:thm:observable-upper}. The converse in \cref{obj:thm:all-estimator-lower} applies to every measurable estimator. Its central ingredient is the weighted approximation and constrained-prior separation result in \cref{obj:thm:weighted-separation-and-frontier}, which supplies the target separation used in the statistical lower bound. \Cref{sec:appendix-weighted,sec:appendix-lower-bound} develop the approximation results and the associated experiment construction.
